\section{Background \& State of the Art}
\label{sec2: background}
\noindent In this section, we provide the biological, computational and mathematical background and context for the 
methods developed in this report. 
We will begin by introducing bacteriophages and the host-prediction problem in Section \ref{sec2.1: host prediction 
problem}, which motivates the need for computational approaches to find a solution. We will then go over the already 
existing and state-of-the-art computational techniques in Section \ref{sec2.2: current approaches}, progressing from the 
classical sequence similarity techniques to the much nore efficient and modern protein embeddings techniques. We will 
also look at the limitations that all the currently exisiting methods have in common, for example the absence of rigourous 
distribution-free uncertainity quantification. \todo{define uncertainty quantification at somepoint in the report}.  In 
the remaining part of the section, \ref{sec2.3: conformal prediction} and \ref{sec2.4: CSA Method} we will introduce the 
mathematical framework of conformal prediction and the score-based aggregation method of in \cite{ochoarivera2024conformal}, 
on which we will build our work.

% However, identifying specific host-range interactions brings some challenges. Traditional laboratory screening via 
% plaque assays is accurate but expensive, labor-intensive, and incapable of scaling to the astronomical number of novel 
% phage genomes generated by modern sequencing \cite{coclet2021global}. Fast, reliable computational prediction tools are 
% therefore essential to prioritize candidate phages prior to experimental validation or to infer host information directly 
% from genomic data.

\subsection{Bacteriophages and Host-Prediction Problem}
\label{sec2.1: host prediction problem}
\noindent Bacteriophages, commonly referred to as phages, are viruses that specifically infect and replicate within 
bacteria. They are one of the most abundant biological entities on Earth with an estimated population of $10^{31}$ in the 
global biosphere \cite{dejonge2019molecular}, and therefore play a very important role in shaping the microbial communities 
and evolution. In the recent years the interest in bacteriophages has increased due to the growing threat of antimicrobial 
resistance and the development of phage therapy as an alternate treatment method to the traditional antibiotics 
\cite{leprince2025strengthening}.

\vspace{0.15cm} 
\noindent One of the key characteristics of of phages is its host range, i.e. the set of bacterial hosts that the phage 
is capable of infecting. The host range of a phage is not a fixed property, it is determined by molecular interactions 
between the phage and the bacterial surface structures during the cycle of infection \cite{dejonge2019molecular}, 
\cite{holtappels2023drivers}. Each phage attaches itself to the bacterial host, injects its viral genetic material which 
are able to replicate inside the host. Therefore, some of the phages only infect a few hosts (have a small host range) as 
they can only infect a limited number of bacterial strains while there are other which are capable of infecting a broad 
range of bacterias. 

\vspace{0.15cm}
\noindent Identifying which phage infects which host, can have several applications both in environmental and clinical 
studies. In clinical, the host range helps us determine how suitable a phage is for designing the treatment for phage 
therapy (an alternate method to antibiotics). In environmental studies it is valuable to know the host being 
infected to understand the ecological interactions between viruses and microbial communities \cite{holtappels2023drivers}
and evolution in the biosphere.

\vspace{0.15cm}
\noindent However, host range has also brings its own obstacles. In phage therapy, we need to find precisely the phage 
that infects the given bacterial pathogen which is the cause of the infection, and for that we need a fast 
reliable host prediction method. Traditionally, the method of determining the host range used to depend on laboratory 
experiments in which candidate phages were tested aginst a number of bacterial hosts; although this method was able to 
give accurate reliable results but it was extremely expensive and time consuming. Given the rapid increase in 
the phage genomes generated by modern technology, the number of possible phage-host interacions has increased 
astronomically, and this leads us to the need of having a compuational prediction method. We need a way which can 
complement the lab experiment by providing with a shorter list of candidate phages before we move on to expermiental 
validation \cite{coclet2021global} or even directly infer the host information directly from a given phage.

\vspace{0.15cm}
\noindent We can formulate this host prediction problem as follows:
% over a probability space $\left(\Omega, \mathcal{F}, \mathbb{P}\right)$. 
let $\mathcal{X} \subseteq \mathbb{R}^d$ be the feature space representing phage representations, and let $\mathcal{H} 
= \{h_1, \cdots, h_N\}$ be the finite label space of possible bacterial hosts. 

\vspace{0.15cm}
\noindent In a single-host classification setting, we observe random pairs $(X, H) \sim \mathbb{P}_{X,H}$, where 
$X \in \mathcal{X}$ is a feature vector representing a phage and $H \in \mathcal{H}$ represents the true host. 
In a general multi-host setting, the target is a non-empty subset of hosts $H^* \subseteq \mathcal{H}$, i.e., 
$H^* \in \mathcal{P}(\mathcal{H}) \setminus \{\emptyset\}$, where $\mathcal{P}(\mathcal{H}) = 2^{\mathcal{H}}$ denotes 
the power set of $\mathcal{H}$ \todo{refer to correct section}. 

\vspace{0.15cm}
\noindent Our primary objective is to construct a set-valued prediction mapping:
$$ C: \mathcal{X} \longrightarrow \mathcal{P}(\mathcal{H})$$
which maps an unlabelled test phage $X_{n+1} \in \mathcal{X}$ to a prediction set $C(X_{n+1}) \subseteq \mathcal{H}$ 
such that it guarantees that the true host is included in the prediction set with a with a user specified statistical 
confidence, $1 - \alpha \in (0, 1)$:
\begin{equation*}
\label{eq: core coverage guarantee objective}
\mathbb{P}\left( H_{n+1} \in C(X_{n+1}) \right) \ge 1 - \alpha
\end{equation*}
\noindent where $\alpha \in (0, 1)$ denotes the target error probability.
